\documentclass[10pt,a4paper]{amsart}
\usepackage[margin=19mm]{geometry}
\usepackage{amsmath,amssymb,mathtools}
\usepackage[scr=rsfs,cal=euler]{mathalfa}
\usepackage{xfrac}
\usepackage[unicode,hidelinks,bookmarks=false]{hyperref}
\newcommand{\ie}{i.e.\ }
\newcommand{\cf}{cf.\ }
\newcommand{\resp}{respectively\ }
\newcommand{\Subsection}{\subsection}
\providecommand{\nobreakdash}{\nobreak\hbox{-}}
\setlength{\emergencystretch}{2em}
\multlinegap0pt
\usepackage{bm}
\usepackage{bbm}
\usepackage{dsfont}
\usepackage{xspace}
\usepackage{cancel}
\usepackage{mathtools}
\usepackage{amsthm}
\makeatletter
\@ifundefined{proposition}{\newtheorem{proposition}{Proposition}}{}
\@ifundefined{theorem}{\newtheorem{theorem}{Theorem}}{}
\makeatother
\title[A new approach to the Poincar\'e center problem]{A new approach to the Poincar\'e center problem}
\author{Isaac A. Garc\'ia, Jaume Gin\'e}
\date{Source: arXiv:2603.09941v3 (CC BY 4.0)}
\begin{document}
\maketitle

\noindent\emph{Source and licence.} A new approach to the Poincar\'e center problem. Version arXiv:2603.09941v3 (CC BY 4.0), distributed under the Creative Commons Attribution 4.0 International licence, as recorded in the arXiv licence metadata for this exact version and shown on the abstract page https://arxiv.org/abs/2603.09941v3. Full rights notice: licenses/arxiv-2603.09941-CC-BY-4.0.txt.\par\smallskip

\noindent\textit{Editorial scope.} Counted unit: the Proposition whose source label is \texttt{prop-seihomo-35} in the article (source0.tex lines 328--336, source label \texttt{prop-seihomo-35}), delivered together with the article's own complete original proof of it (source0.tex lines 338--360, one \texttt{proof} environment of 23 lines). No mathematics is written, repaired or re-derived: statement and proof are the article's own text line for line. Required context retained uncounted: V-Laurent (lines 86--88); teo-Pi-Algorithm (lines 178--184); edo-para-vj (lines 221--223); vm-explicit (lines 229--231); semi-mM (lines 318--320). Every definition, display and label of those lines is retained unchanged. Omitted from this package and not redistributed: 1--85, 89--177, 185--220, 224--228, 232--317, 321--327, 337--337, 361--592. Nothing inside a retained excerpt is omitted or altered: every retained body is delivered exactly as the article's lines are written, so no omission receipt of any kind accompanies this package. Citations: the retained excerpts carry no citation command at all, so no bibliography entry of the article is transcribed and no reference is redistributed. Reference adaptation: none was needed and none was made; every reference inside the delivered unit resolves inside it, so no unresolved reference or citation is produced. Printed slips kept literal: no printed word, symbol, hypothesis or number of the counted statement or of its proof was altered. Layout and class: the article's own class is \texttt{amsart} with packages including color, amsmath, amsthm, amssymb, amscd, amsfonts, amsbsy, cite, epsfig, afterpage, paralist, psfrag, mathrsfs, verbatim; the wrapper is the lane's pinned \texttt{amsart} editorial wrapper at 10pt on A4, which restates the theorem-like environments the retained text needs and adds the article's own macro definitions that the retained text needs (none), with extra wrapper packages (bm, bbm, dsfont, xspace, cancel, mathtools). The delivered numbering is the unit's own. Assets: no retained span contains a figure environment, a graphics-inclusion command, a PDF-page inclusion, a table or a TikZ picture; no image file is copied into this package, the package declares no asset dependency and no image file is redistributed with it. Licence: version arXiv:2603.09941v3 of this article is distributed under the Creative Commons Attribution 4.0 International licence, as recorded in the arXiv licence metadata for this exact version and shown on the abstract page https://arxiv.org/abs/2603.09941v3. A full rights notice is delivered at licenses/arxiv-2603.09941-CC-BY-4.0.txt. Characters: every retained excerpt is pure ASCII, so the delivered engine, XeLaTeX with its default Latin Modern text font, reports no missing character.
\begin{equation}\label{V-Laurent}
V(\varphi, \rho) = \sum_{j \in \mathbb{Z}} v_j(\varphi) \rho^{j},
\end{equation}

\begin{theorem} \label{teo-Pi-Algorithm}
We restrict to the parameter space $\Lambda \setminus \Lambda_{pq}$ and, by a rotation, we take $0 \not\in \Omega_{pq}$. Then the Laurent expansion \eqref{V-Laurent} with initial condition $V(0, \rho)$ analytic at $\rho=0$ becomes a convergent Taylor expansion
\begin{equation}\label{Taylor-V}
V(\varphi, \rho) \Big|_{I_{pq}} = \sum_{j \ge m} v_j(\varphi) \, \rho^j,
\end{equation}
with $m \geq 1$ and $I_{pq} = [0, 2 \pi] \setminus \Omega_{pq}$.
\end{theorem}

\begin{equation}\label{edo-para-vj}
\mathcal{L}^{a}_j(v_j) + Q_j = 0
\end{equation}

\begin{equation}\label{vm-explicit}
v_m(\varphi; C_m) = C_m G_0(\varphi) \exp\left( (1-m) \, \mathcal{P}(\varphi) \right)
\end{equation}

\begin{equation}\label{semi-mM}
\dot{x} = P_n(x,y) = -y^n + \sum_{\stackrel{i+j=n}{\scriptscriptstyle i \neq 0}} a_{ij} x^i y^j, \ \ \dot{y} = Q_N(x,y) = x^N + \sum_{\stackrel{i+j=N}{\scriptscriptstyle j \neq 0}} b_{ij} x^i y^j,
\end{equation}

\begin{proposition}\label{prop-seihomo-35}
Let the origin of the monodromic $(3, 5)$-semi-homogeneous family \eqref{semi-mM} be a center. Then $L_1 = 0$, where
\begin{eqnarray*}
L_1 &=& a_{12}^5 + 5 a_{12}^3 a_{21} + 5 a_{12} a_{21}^2 - 2 a_{21}^5 b_{05} + a_{12} a_{21}^4 b_{14} -  a_{12}^2 a_{21}^3 b_{23} \\
 &  & - 2 a_{21}^4 b_{23} + a_{12}^3 a_{21}^2 b_{32} + 3 a_{12} a_{21}^3 b_{32} -
 a_{12}^4 a_{21} b_{41} - 4 a_{12}^2 a_{21}^2 b_{41} -
 2 a_{21}^3 b_{41}.
\end{eqnarray*}
\end{proposition}

\begin{proof}
The $(3, 5)$-semi-homogeneous family \eqref{semi-mM} has $W(\mathbf{N}(\mathcal{X})) = \{ (1, 2), (1, 1) \}$. Taking polar coordinates we have that $0 \in \Omega_{11}$, hence we permute coordinates $(x,y) \mapsto (y,x)$ and, taking again polar coordinates in the transformed system, we get that $G_0(\varphi) = \cos^2\varphi (- \cos^2\varphi + a_{12} \sin\varphi \cos\varphi  + a_{21} \sin^2\varphi) \leq 0$ and $\Omega_{11} = \{\pi/2, 3 \pi/2\}$ when we restrict to $\Lambda$.

The equation $\Theta(\varphi, \rho) = G_0(\varphi) + G_2(\varphi) \rho^2 = 0$ is a quadratic equation in $\rho$. A necessary condition to have real positive solutions $\rho=\rho^*(\varphi, a, b) = \sqrt{-G_0/G_2}$ in a neighborhood of $\Omega_{11}$ and $\Lambda$ is that $G_2 > 0$ there. But $G_2$ is continuous and $G_2|_{\Omega_{11}} = -1 < 0$, hence $\Lambda_{11} = \emptyset$.

We have that $R_1(\varphi)/G_0(\varphi) = \tan\varphi$ with primitive $p(\varphi) = -\log(|\cos\varphi|)$ noncontinuous at $\Omega_{11}$ going to $+ \infty$. Then we cannot build any primitive bounded at $\Omega_{11}$. Then the function $v_m$ with $m>1$ given in \eqref{vm-explicit} is $v_m(\varphi; C_m) = C_m G_0(\varphi) \, |\cos^{m-1}\varphi|$ with $C_m \neq 0$.

Instead of analyzing the next coefficient $v_{m+1}$, we prefer to take $m=1$ that is justified using Theorem \ref{teo-Pi-Algorithm} with any initial analytic condition $V(0,\rho) = c_1 \rho + \cdots$ with $c_1 \neq 0$. Then
$v_1(\varphi; 1) = G_0(\varphi)$ with $C_1 =1$ without loss of generality.


The next coefficient satisfies an homogeneous equation $\mathcal{L}^{a}_2(v_2) = 0$ with general solution $v_2(\varphi; C_2) = C_2 G_0(\varphi) \, |\cos\varphi|$. We take $C_2=0$ (hence just $v_2\equiv 0$) choosing the compatible analytic initial condition $V(0,\rho) = c_1 \rho + O(\rho^3)$ so that $v_2(0; C_2) = C_2 = 0$.

The differential equation \eqref{edo-para-vj} with $j=3$ that $v_3$ satisfies is no longer homogeneous since $Q_3 \not\equiv 0$. We let $v_3^h(\varphi)$ to be the solution of the associated homogeneous equation $\mathcal{L}^{a}_3(v_3^h) = 0$. Up to a multiplicative constant it has the form $v_3^h(\varphi) = -2 \cos^2\varphi  G_0(\varphi)$. We express $v_3(\varphi; C_3) = v_3^h(\varphi) (w_3(\varphi)+ C_3)$ and we obtain that
$$
w_3(\varphi) = \frac{L_1}{\sqrt{-\Delta}} \, F_3(\varphi) + G_{3}(\varphi)
$$
being $\Delta = a_{12}^2 + 4 a_{21} < 0$ in $\Lambda$ and $F_3(\varphi)$ is a continuous function on $[0, 2 \pi]$ constructed piecewise where in each of the 3 components of $\Omega_{11} \setminus [0, 2 \pi]$ it is just to add some constant to the function
$$
{\rm arctan}\left( \frac{a_{12} -2 a_{21} \tan\varphi}{\sqrt{-\Delta}}\right).
$$
Moreover, $G_{3}$ is a function defined in $\mathbb{S}^1 \setminus \Omega_{11}$. However, $\lim_{\varphi \to \varphi^*} v_3^h(\varphi) G_{3}(\varphi)= -1/2$ when $\varphi^* \in \Omega_{11}$ and $v_3^h \, G_{3}$ can be extended by continuity to $\Omega_{11}$. Clearly this property is also shared by $v_3$. The necessary center condition determined by imposing $v_{3}$ be defined in $\mathbb{S}^1$ turns out to be that $L_1=0$ since $F_3$ is monotonous, hence not $2 \pi$-periodic.
\end{proof}

\end{document}
